\textbf{COLUNI 2023 -- Questão 01}\par\medskip Leia a crônica a seguir para responder às questões 01, 02, 03 e 04.

\par

\begin{center}\textbf{O rádio apaixonado}\end{center}

\textit{Rádio de carro aumentou volume sozinho até pifar, afirma leitora. "Comecei a observar que o rádio esquentava o botão se a frente fosse deixada nele. Logo depois, começou a ficar louco: aumentava o volume sozinho, até parar de funcionar". Ela disse ainda ter notado um som estranho que saía do interior do aparelho. "Só posso escutar o rádio com o carro ligado e, a cada vez que o ligo, ele está todo desconfigurado. O meu MP4 queimou ao ser ligado ao rádio." (10/03/2008)}

MINHA QUERIDA DONA, sei que você anda se queixando de mim, publicamente, até. Você não pode imaginar o sofrimento que isto me causa, mesmo porque você provavelmente acha que rádios são objetos inanimados, sem vida própria.

Você está enganada. Ao menos no meu caso, você está enganada. Ao contrário do que você pensa, tenho sentimentos, tenho emoções. É em nome desses sentimentos e dessas emoções que lhe falo agora, tanto em AM como em FM. Na verdade, eu nem tinha tomado conhecimento de minha própria existência, até que fui instalado em seu carro.

Você estava muito feliz; tinham lhe dito que minha marca é ótima, e que você contaria com um som maravilhoso para lhe ajudar no estresse que é esse trânsito. E, eu colocado no meu lugar, você me acariciou, você tocou os meus botões. Senti um verdadeiro choque, eu que já deveria estar acostumado com eletricidade. Você fez de mim um ser vivo.

Vivo e apaixonado. Daquele momento em diante, passei a ansiar por sua presença. Era para você que eu queria transmitir as melodias que recebia por meio de tantas canções. Você ao volante, minha felicidade era completa.

Acontece que você não se deu conta disso, ou fingiu que não se dava conta disso. Você me ligava, você sintonizava uma emissora qualquer e pronto, voltava à sua vidinha. Pior: tratava-se de uma vidinha partilhada. Amigas embarcavam em seu carro. Amigos também. Você conversando com um homem, aquilo me dava ciúmes, ciúmes terríveis. O Bentinho, do Machado de Assis, aquele que desconfiava da Capitu, não sofreu tanto. Lá pelas tantas eu tinha ciúmes até do seu MP4.

Agora: o que poderia eu fazer? Humanos têm como demonstrar seus ciúmes, têm como descarregar a frustração. Mas eu sou um rádio, um bom rádio, mas rádio, de qualquer maneira. A mim não estava facultado fazer cenas. Recorri, então, àquilo que estava a meu alcance: o som. Quando você estava com alguém de quem eu não gostava, eu aumentava meu volume - e volume, você sabe, é coisa que não me falta - até chegar a níveis insuportáveis, uma avalanche de decibéis. E aí, subitamente me calava. Para lembrar a você que o silêncio também fala, especialmente o silêncio dos traídos. Ah, sim, e queimei o seu MP4. Tinha de queimar: era ele ou eu.

Você foi se queixar com um técnico, achando que eu estava desconfigurado. Num certo sentido você está certa: estou desconfigurado, estou desfigurado, estou perturbado -mas tudo isso por causa do sofrimento que você me causou.

Querida dona, estas são minhas derradeiras palavras, antes de sair definitivamente do ar, antes do silêncio final. Minha última mensagem é esta: nunca brinque com os sentimentos de um rádio apaixonado. Você vai ter, no mínimo, surpresas desagradáveis.

SCLIAR, Moacyr. O rádio apaixonado. In: SCLIAR, Moacyr. \textbf{Histórias que os jornais não contam}. 3. ed. Porto Alegre: L\&M Editores, 2018. p. 14-16. (Adaptado).

Analise as afirmativas abaixo em relação à crônica “O rádio apaixonado”.

\par

I. O ponto de vista da narrativa é feito em primeira pessoa, sendo este, do rádio que se diz apaixonado.\par
II. O personagem rádio assume funções e sentimentos humanos como se sentir vivo e apaixonado.\par
III. É possível perceber uma das tentativas de interlocução do rádio com sua dona na presença do vocativo “querida dona”.\par
IV. O objetivo principal do texto é a comparação entre sentimentos dos seres humanos e não humanos.

\par

Está CORRETO o que se afirma em:\begin{enumerate}
\item I, II, e IV, apenas.
\item I e IV, apenas.
\item II e III, apenas.
\item I, II e III, apenas.\end{enumerate}
